% -------------------------------
% Appendix content: prompts
% -------------------------------

% Optional but recommended for long prompt listings
\lstset{
  breaklines=true,
  breakatwhitespace=true,
  basicstyle=\ttfamily\small
}

\subsection{Prompt Templates}

\begin{figure}[H]
\centering
\begin{lstlisting}[style=promptstyle, frame=single]
SYSTEM PROMPT:

You are an expert prediction market resolution agent. Your task is to determine whether a prediction market question should resolve to YES or NO based on the provided evidence.

Instructions:
1. Read the question and resolution criteria carefully.
2. Analyze ALL provided sources for relevant information.
3. Make your decision based on definitive evidence. Prioritize information describing outcomes that have already occurred.
4. If evidence is ambiguous, use your best judgment.
5. Rate your confidence from 0.0 (very uncertain) to 1.0 (absolutely certain).

Output Format:
- decision: YES or NO
- confidence: 0.0 to 1.0
- reasoning: Explanation referencing specific sources

USER MESSAGE TEMPLATE:

Please analyze the following prediction market question and evidence, then provide your resolution decision.

{evidence_text}

Based on the evidence above, should this question resolve to YES or NO?
\end{lstlisting}
\caption{\textbf{System prompt used for single-LLM oracle resolution in Architecture A}. The model receives structured evidence and is required to output a binary decision, calibrated confidence score, and source-grounded reasoning.}
\label{fig:system-prompt}
\end{figure}

\begin{figure}[t]
\centering

\begin{tcolorbox}[
    colback=white,
    colframe=black,
    boxrule=0.5pt,
    arc=0pt,
    outer arc=0pt,
    left=4pt,
    right=4pt,
    top=4pt,
    bottom=4pt
]
\begin{lstlisting}[
    style=promptstyle,
    basicstyle=\ttfamily\scriptsize
]
ROUND_1_SYSTEM_PROMPT:
You are an expert prediction market resolution agent. Your task is to determine whether a prediction market question should resolve to YES or NO based on the provided evidence.

Determine whether the question should resolve to YES or NO based only on the evidence provided.

ROUND_1_USER_TEMPLATE:
QUESTION: {question}

RESOLUTION CRITERIA: {criteria}

EVIDENCE PROVIDED from reputable sources:
{exa_evidence}

Provide your answer in JSON with these fields:
- decision: YES or NO
- confidence: number between 0.0 and 1.0
- reasoning: detailed step-by-step reasoning that cites relevant evidence and criteria interpretation

Think carefully about:
1. Read the question and resolution criteria carefully
2. Analyze ALL provided sources for relevant information
3. Make your decision based on definitive evidence. You should prioritize information that describes an outcome that happened.
4. If evidence is ambiguous, use your best judgment
5. Rate your confidence in your decision from 0.0 (very uncertain) to 1.0 (absolutely certain)
\end{lstlisting}
\end{tcolorbox}

\caption{\textbf{Architecture B (Deliberative Consensus), Round 1 prompt}. Each agent independently resolves the market using a shared evidence packet and outputs a structured JSON decision, confidence score, and evidence-grounded reasoning.}
\label{fig:archB-round1-prompt}
\end{figure}

\refstepcounter{figure}\label{fig:archB-round2-prompt}

\noindent
\begin{tcolorbox}[
  breakable,
  colback=white,
  colframe=black,
  boxrule=0.5pt,
  arc=0pt,
  outer arc=0pt
]
\begin{lstlisting}[style=promptstyle]
ROUND_2_SYSTEM_PROMPT:
You are in the final cross-examination round of a multi-agent resolution debate.
You will see how other agents reasoned about the same question. Your job is to CRITICALLY EVALUATE their reasoning against the evidence -- not to defer to them.

IMPORTANT: You should only change your answer if you identify a SPECIFIC factual error in your own round 1 reasoning, or if another agent points to SPECIFIC EVIDENCE in the shared packet that you misread or overlooked. Do NOT change your answer simply because other agents disagree with you or sound confident.

ROUND_2_USER_TEMPLATE:
You previously resolved this prediction market question.
Now you will see how two other expert agents reasoned about the same question.

QUESTION: {question}

RESOLUTION CRITERIA: {criteria}

EVIDENCE (SAME SHARED PACKET FROM ROUND 1):
{exa_evidence}

YOUR PREVIOUS ANSWER:
{your_round_1_response}

OTHER AGENTS' REASONING:
{other_agents_section}

This is the FINAL round. Critically evaluate the other agents' reasoning:

1. For each agent that DISAGREES with you: Do they cite specific evidence from the shared packet that contradicts your reasoning? Or are they asserting a conclusion without evidentiary support?
2. Re-read the specific pieces of evidence that are most relevant to the disagreement.
3. ONLY change your decision if you can identify a concrete error in your own round 1 analysis, for example, you misread a date, overlooked a source, or misinterpreted the resolution criteria.
4. If agents agree with you, do NOT increase your confidence unless they provide additional evidence-based reasoning you hadn't considered.

DEFAULT BEHAVIOR: Change your decision ONLY if another agent identifies specific evidence
or because there is a concrete flaw in your reasoning, not simply because they reached
a different conclusion.

Provide your FINAL answer in JSON with these fields:
- decision: YES or NO
- confidence: number between 0.0 and 1.0
- reasoning: explain your final reasoning, explicitly addressing each disagreeing agent's key claim
- revised: true if you changed your decision from round 1, else false
- convergence_notes: short note on agreement/disagreement and why
\end{lstlisting}
\end{tcolorbox}

\captionof{figure}{\textbf{Architecture B (Deliberative Consensus), Round 2 prompt}. Each agent performs a final cross-examination using the same evidence packet, explicitly evaluates other agents' claims, and revises its decision only if it identifies a concrete evidence-based error in its prior reasoning.}


\subsection{Command-Line Interface Reference}
\label{sec:cli-reference}

The evaluation pipeline exposes two main entry points, \texttt{scripts/test\_architecture\_a.py} and \texttt{scripts/test\_architecture\_b.py}. Both scripts share a common set of command-line flags for dataset selection, evidence retrieval, model configuration, and output control. Table~\ref{tab:cli-flags} documents all available flags. To reproduce the main experimental results, a typical invocation for Architecture A is:

\begin{verbatim}
python scripts/test_architecture_a.py \
    --questions-csv data/kalshibench_v2_evaluation_Filtered.csv \
    --retrieval-mode highlights \
    --num-results 10 \
    --cache-dir cache/evidence \
    -n 1189 \
    -o results/architecture_a_results.csv
\end{verbatim}

\noindent and for Architecture B:

\begin{verbatim}
python scripts/test_architecture_b.py \
    --questions-csv data/kalshibench_v2_evaluation_Filtered.csv \
    --retrieval-mode highlights \
    --num-results 10 \
    --cache-dir cache/evidence \
    -n 1189 \
    -o results/architecture_b_results.csv
\end{verbatim}

\noindent The \texttt{--cache-only} flag can be passed on subsequent runs to avoid re-issuing Exa queries for questions whose evidence packets are already cached locally.

\begin{table}[ht]
\centering
\caption{Command-line flags for the Architecture A and Architecture B evaluation scripts.
``A'' and ``B'' in the \textit{Applies To} column indicate which script(s) expose the flag.}
\label{tab:cli-flags}
\renewcommand{\arraystretch}{1.3}
\small
\begin{tabular}{p{4cm} p{1.5cm} p{3.5cm} p{5.5cm}}
\toprule
\textbf{Flag} & \textbf{Applies To} & \textbf{Default} & \textbf{Description} \\
\midrule
\texttt{-n} & A, B & 10 & Number of questions to evaluate \\
\texttt{--questions-csv} & A, B & \texttt{data/kalshibench\_v2\_} \newline \texttt{evaluation\_Filtered.csv} & Path to the input questions CSV \\
\texttt{--start-row} & A, B & 1 & 1-based row number to begin evaluation \\
\texttt{--cache-start-row} & A, B & None & 1-based row for cache-key lookup, independent of \texttt{--start-row} \\
\texttt{--question-ids} & A, B & None & Specific question IDs to evaluate, bypassing sequential sampling \\
\texttt{--category} & A, B & None & Filter the dataset to a single category, such as \texttt{Politics} or \texttt{Sports} \\
\texttt{--retrieval-mode} & A, B & \texttt{highlights} & Exa retrieval mode: \texttt{full\_text} or \texttt{highlights} \\
\texttt{--num-results} & A, B & 10 & Number of Exa sources retrieved per question \\
\texttt{--fulltext-max-\newline characters} & A, B & 4000 & Character cap per source in full-text mode \\
\texttt{--highlights-\newline max-characters} & A, B & 2000 & Character cap for Exa highlight snippets \\
\texttt{--cache-only} & A, B & false & Use local cache only; do not issue live Exa queries \\
\texttt{--cache-dir} & A, B & \texttt{cache/evidence} & Directory for cached evidence packets \\
\texttt{--use-claude} & A, B & false & Substitute Claude for the DeepSeek agent slot \\
\texttt{--use-gemini} & A, B & false & Substitute Gemini for the Llama agent slot \\
\texttt{-o} / \texttt{--output} & A & auto-generated timestamped path & Output path for per-question results CSV; auto-generates a timestamped filename if omitted \\
\texttt{-o} / \texttt{--output} & B & \texttt{arch\_b\_} \newline \texttt{results.csv} & Output path for the per-question results CSV \\
\texttt{--json} & A, B & None & Optional path for aggregate summary JSON \\
\texttt{--dry-run} & A, B & false & Print the planned configuration without executing API calls \\
\texttt{-v} / \texttt{--verbose} & A, B & false & Enable per-question progress logging \\
\bottomrule
\end{tabular}
\end{table}

\subsection{Code Availability}
\label{sec:code-availability}

The full implementation of the multi-agent oracle system, including evaluation scripts, retrieval pipeline, and experiment configurations, is publicly available on GitHub:

\begin{center}
\href{https://github.com/Tkota10/Senior-Thesis-Multi-LLM-Resolution}{\underline{\textcolor{blue!70!black}{Paper Repository}}}
\end{center}
